\documentclass[12pt, letterpaper]{article}
\usepackage[T1]{fontenc}
\usepackage[spanish]{babel}
\usepackage[margin=3.5cm]{geometry}
\usepackage{setspace}
\usepackage{charter}
\usepackage{float}
\usepackage[autostyle=true]{csquotes}


% Paquetes matemáticos y gráficos (útiles si el ensayo lo requiere)
\usepackage{amsmath}
\usepackage{graphicx}
\usepackage{hyperref}
\usepackage{bookmark}

\title{LTII: Visión y espectativas sobre el proyecto final}
\author{@Char[50]}
\date{\today}
% Cambio de la cardinalida para la secciones 
\renewcommand{\thesection}{\Roman{section}}
\renewcommand{\thesubsection}{\alph{subsection}}


\begin{document}
\maketitle


\newpage
\section{Meta del proyecto}

\vspace{0.5cm}

Mi intensión con este trabajo es ilustrar al espectador sobre \enquote{lo bueno y lo malo} que la plataforma de \textbf{Steam} carga consigo. Puesto que, apesar de todo lo bonito que pueda ser, sería un locura negar el enorme poder que tiene sobre el ámbito del \textit{Gaming} en \textit{PC}. Siendo una plataforma que dictámina \enquote{el ritmo del juego} y que cada decisión es plenamente libre de consulta a terceros, considero bastante sensato tomarlo en cuenta.

\vspace{0.5cm}

Por otro lado, también es cierto que cómo plataforma aporta mucho. Steam ha sido reconocida durante años como \enquote{pro-consumidor} y \enquote{pro-usuarios} poniendo, en gran medida, a sus usarios por delante de cualquier otra cosa. Lo que naturalmente a creado bastante comunidad y le ha otorgado la preferencia y popularidad que ostenta. 

\vspace{0.5cm}

Un ejemplo bastante claro de lo que quiero contar, puede verse reflejado en la situación que enfrenta casí cualquier proyecto \textit{Indie} que sale al mercado hoy en día. Steam ofrece a los \textit{publishers} tanto una basta comunidad como varias herramientas para lanzar y promocionar el juego pero al mismo tiempo involucra 2 situaciones que son imposibles de ignorar: 
\begin{itemize}
    \item Existe una comisión del 30\% que Steam cobra por manejar el lanzamiento. \textit{\enquote{Transacción que no siempre es del todo viable}}
    \item De elegir otra plataforma u otro medio para hacerlo, la cantidad de foco mediático se reduce inmensamente. \textit{\enquote{Casi una sentencia de muerte, sobre todo para un lanzamiento de esta categoría}}
\end{itemize}

\vspace{0.5cm}

\section{Planteamiento general}

\vspace{0.5cm}

Para lograr compaginar la narrativa que estoy desarrollando y el uso de datos cualitativos referentes a \enquote{ciertos} tópicos que me ayuden a representar \enquote{tanto lo bueno como lo malo} de Steam, el proyecto se apoyará en la técnica del \textit{Scrolltelling}. De tal manera, que el espectador sea capaz de avanzar a través del texto mientras algunos tipos de gráficas acompañan el punto principal de cada sección. La idea detrás de usar dicha técnica es aprovechar los datos mediante un recurso visual llamativo y consistente, para potenciar cada punto. 

\newpage
\section{Disposición de elementos}
Para no restarle protagonismo ni a la narrativa ni a las gráficas, he tomado la desición distribuir la información como se muestra en el siguiente boceto: 


\end{document}